\documentstyle{article}

\setlength{\topmargin}{-1.4cm}
\setlength{\marginparwidth}{2cm}
\setlength{\oddsidemargin}{0.3cm}
\setlength{\evensidemargin}{0.3cm}
\setlength{\textwidth}{15.8cm}
\setlength{\textheight}{24cm}
\newcommand{\sub}[1]{_{\rm #1}}

\begin{document}
\begin{center}
\vspace*{3.0 cm}
{\LARGE \bf SimLine - Radiative transfer in molecular lines\\}
\vspace{1.2cm}
{\Large \bf Version 2.18 (03/22/2026)\\}
\vspace{0.8cm}
{\large  Dr.~V.~Ossenkopf-Okada\\\vspace{0.2cm}
I. Physikalisches Institut der Universit\"at zu K\"oln\\}
\vspace{3.8cm}
{\LARGE \bf Users Manual\\}
\vspace{5cm}
\end{center}

\begin{abstract}
{\normalsize SimLine is a FORTRAN code to compute the profiles of
molecular transition lines in spherically symmetric clouds with 
arbitrary density, temperature and velocity structure.
Turbulence and clumping effects are treated in a local statistical 
approximation combined with a radial dependence of the correlation
parameters. The code consists of two parts:
the self-consistent solution of the balance equations for all level
populations and energy densities at all radial points
and the computation of the emergent line profiles observed
from a telescope with finite beam width and arbitrary offset. 
All numerical discretizations are done in an adaptive way,
the user has full error control and the line profiles may
be computed in almost arbitrary accuracy. An accelerated lambda
iteration is used to solve the equation system. The optical
depths in the lines may vary between about -5 and +5000.
The input of physical cloud parameters, telescope parameters, and 
numerical control parameters may be done either interactively 
or as file input. Molecular data are read from table files.}
\end{abstract}
\thispagestyle{empty}
\newpage

\section{The basic theory}
\subsection{The solution of the radiative transfer problem}

The general formulation of the frequency dependent radiative transfer 
equation
\begin{equation}
{dI(\nu,\vec{r},\vec{n})\over ds} = - \kappa(\nu,\vec{r},\vec{n})
I(\nu,\vec{r},\vec{n}) + \epsilon(\nu,\vec{r},\vec{n})
\end{equation}
may be reduced in the spherically symmetric case down to three
independent coordinates. It is sufficient to consider radiation
in the direction of the $z$ axis and to neglect the $\phi$ dependence
of the radius vector. Then, the radiative transfer equation can be written 
as
\begin{equation}
{dI_z(\nu,p,z)\over dz} = -\kappa(\nu,p,z) I_z(\nu,p,z)
+ \epsilon(\nu,p,z)
\end{equation}
where $p$ denotes the displacement variable ($p=\sqrt{x^2+y^2}$).
The $z$ axis is chosen in the direction towards the observer.
Since the spatial grid is adjusted at distinct radial points,
the program does not use the form of the radiative transfer
equation with source function and derivative on the optical
depth, but with emission and absorption coefficients.

For the numerical solution of the radiative transfer equation
it is convenient to use stepwise the solution of the integral
equation allowing a reduction of the error to third order
\begin{equation}
I_z(\nu,p,z_i) = \exp \left( -\int_{z_{i-1}}^{z_i} \!\!\!
\kappa(\nu,p,z) dz \right) \left[ I_z(\nu,p,z_{i-1}) +
\int_{z_{i-1}}^{z_i} \!\!\!\epsilon(\nu,p,z) \exp \left( \int_{z_{i-1}}^z
\kappa(\nu,p,z') dz' \right) dz \right]
\label{transfer}
\end{equation}
The incident radiation at the outer boundary of the cloud is
assumed to follow a black body spectrum with temperature $T_{bg}$.
\begin{equation}
I\sub{bg}(\nu) = {2 h \nu^3 \over c^2}\;\left[\exp \left(h \nu \over
k \times T_{bg}\right) -1\right]^{-1}
\end{equation}
For the cosmic background radiation $T_{bg}=2.73K$.

The program assumes complete redistribution of energy between
the molecules at a given level, so that one can define a unique set
of level populations $n_j$ at any point $\vec{r}$ which is independent
from frequency or the actual velocity of a molecule. Furthermore,
this means that the line profiles for emission and absorption
are the same.  In molecular clouds it is sufficient to consider
pure Doppler broadening so that the profiles are given by 
\begin{equation}
\Phi(\nu)={1 \over \sqrt{\pi} \sigma} \exp{}\left(-{(\nu-\nu_0(1+v_z/c))^2
\over \sigma^2}\right)
\label{phi}
\end{equation}
where $v_z$ is the systematic velocity in $z$ direction, $\sigma$ is the
local line width, and $\nu$ is the 
frequency in the observers frame. (Without loss of generality, the program 
assumes that the centre of the cloud is in rest within the observers frame.).

With the assumtion of complete redistribution, the absorption and 
emission coefficients for a given transition are determined by
\begin{eqnarray}
\kappa_{j,l}(\nu)&=&{h\nu_0 \over c} \;(n_l B_{l,j}-n_j B_{j,l}) \Phi(\nu)\\
\epsilon_{j,l}(\nu)&=&{h\nu_0 \over 4\pi}\; n_{j} A_{j,l} \Phi(\nu)
\label{kappa}
\end{eqnarray}
where stimulated emission is treated as negative absorption.

The level populations at a given radius $r$ may be computed
from the linear system of balance equations
\begin{equation}
n_j \sum_{l\ne j}\left( A_{j,l}+B_{j,l}u_{j,l}+ c_{j,l}\right) =
\sum_{l \ne j} n_l \left(A_{l,j}+B_{l,j}u_{j,l}+ c_{l,j}\right)
\label{balance}
\end{equation}
where the radiative energy density $u_{j,l}$ for the transition between
the levels $j$ and $l$ at a given point is computed from the integral
\begin{equation}
u_{j,l}(p,z) = {1 \over c} \int_0^{2\pi}d\phi \int_0^\pi d\theta\; {\rm sin} \theta
\int_{-\infty}^{\infty} d\nu \,\Phi_{\phi,\theta}(\nu,\vec{v}) 
I(\nu,p,z,\phi,\theta)\;.
\end{equation}
Exploiting the spherical symmetry, this may be transformed into
\begin{equation}
u_{j,l}(r) = {2\pi\over r c} \int_{-r}^r dz' \int_{-\infty}^\infty
\Phi(\nu) I_z(\nu,p'=\sqrt{r^2-z'^2},z') d\nu
\label{uzint}
\end{equation}
where the radius $r$ is given as $r=\sqrt{p^2+z^2}$.
Due to the symmetry, only radiation parallel to the $z$ axis is 
encountered. For the central line frequency $\nu_0$ within the 
expression for $\Phi$ (Eq. \ref{phi}),
the frequency of the transition $j \rightarrow l$ has to be taken
when computing $u_{j,l}$.

The equation system is truncated whenever the excitation of a level
falls below the chosen accuracy limit or at the maximum level for
which the collision rates are known.

\subsection{The local turbulence approximation}

The turbulence description uses two parameter for each point $\vec{r}$:
the total width of the velocity distribution $\sigma$ giving the
local profile for optically thin lines and the correlation length
$r\sub{corr}$.

The width of the velocity distribution $\sigma$ is composed of a 
turbulent and a thermal contribution
\begin{equation}
\sigma = {\nu_0 \over c} \,\sqrt{ {2kT\sub{kin} \over m} + {2 \over 3} 
\langle v\sub{turb}^2 \rangle}
\end{equation}
where a Maxwellian distribution of turbulent velocities is assumed. 
In the data input, the turbulent velocities are given as FWHM of the 
velocity distribution (${\rm FWHM}(v\sub{turb})=\sqrt{8/3\times{\rm ln}2\;
\langle v\sub{turb}^2 \rangle}$). The thermal velocities are automatically
computed from the kinetic gas temperature.

The long range correlation of the turbulence spectrum as described by
a Kolmogorov exponent can be simulated in 1D by a radial dependence
of the turbulent velocity dispersion $\langle v\sub{turb}^2\rangle 
\propto r^\gamma$. Observations of Fuller \& Myers (1992) suggest 
exponents around $\gamma \approx 0.7$. Several other authors proposed
somewhat smaller values around $\gamma \approx 0.5$.

For the local treatment of clumping in real space or velocity space the
considered volume element is subdivided into numerous
clumps whith a thermal internal velocity dispersion.
From a Gaussian distribution of the density of molecules with a certain 
velocity within the single fragments
\begin{equation}
n(r)=n_0\times\exp(-r^2/r\sub{cl}^2)
\end{equation}
we get an effective absorption coefficient for the whole medium
given by
\begin{equation}
\kappa\sub{eff}=n\sub{cl} \times \pi r\sub{cl}^2 \int_0^{\tau\sub{cl}}
{1-\exp(-\tau) \over \tau} d\tau
\end{equation}
(Martin et al. 1984),
where $n\sub{cl}$ is the number density of clumps which contribute
at the considered frequency and $\tau\sub{cl}=\sqrt{\pi}\,\kappa r\sub{cl}$ is their central opacity.

Assuming a Maxwellian turbulent velocity distribution of clumps
and a thermal velocity dispersion $\sigma\sub{th}$ giving at most
1/3 of the total velocity dispersion $\sigma$, we obtain
an effective absorption coefficient
\begin{equation}
\kappa\sub{eff}(\nu)=n\sub{ges} \pi r\sub{cl}^2 \times A(\tau\sub{cl})
\times {\sigma\sub{th} \over \sigma} \exp{}\left(-{(\nu-\nu_0)^2
\over \sigma^2}\right)
\end{equation}
with 
\begin{equation}
A(\tau)={1\over \sqrt{\pi}}\int_{-\infty}^{\infty}dv
\int_0^{\tau \exp(-v^2)} {1-\exp(-\tau') \over \tau'} d\tau'
\end{equation}
Here, $n\sub{ges}$ is the total number of clumps. In case of
pure turbulence in velocity space, $n\sub{ges}$ is given by the
volume of the single clumps and we obtain
\begin{equation}
\kappa\sub{eff}(\nu)=\kappa(\nu)\times A(\tau\sub{cl})/\tau\sub{cl}\;.
\end{equation}
Consequently, $A(\tau\sub{cl})/\tau\sub{cl}$ is a measure of the reduction
of the opacity due to turbulence. For small clump sizes, $A(\tau\sub{cl})
= \tau\sub{cl}$ so that we are in the microturbulent limit. For
$\tau\sub{cl}\gg 1$ the macroturbulent limit is reached.
In case of real clumping in space, $\kappa\sub{eff}(\nu)$ is further reduced
by the filling factor. In the program, this is simulated by a corresponding
artificial reduction of the molecular abundance.

The clumps size $r\sub{cl}$ is determined by the correlation length
of the velocity or density structure. $r\sub{cl}$ is the length on 
which the abundance of molecules within the same thermal velocity profile 
is reduced by the factor 1/e. Consequently, it can be computed from the 
turbulence correlation length by $r\sub{cl}=r\sub{corr}\times \sigma\sub{th}
/\sigma$.  

One has to keep in mind that this is a statistical approximation,
always assuming a complete Maxwellian velocity distribution within each
volume element even if the single clumps become relatively large.
Furthermore, the source function $S=\epsilon(\nu)/\kappa(\nu)$ is not 
influenced in this local approach.

\subsection{The central H~II region}

For the simulation of the internal heating produced by a central continuum 
source in the cloud, it is possible to assume an H~II region in the cloud core.
The H~II region is characterized by two parameters, the electron density
$N\sub{e}$ and the kinetic electron temperature $T\sub{e}$.

The absorption coefficient for electron-ion bremsstrahlung in the
Rayleigh-Jeans approximation is given by:
\begin{equation}
\kappa(\nu)={8 \over 3\sqrt{2\pi}}\; {e^6 \over (4\pi\epsilon_0 m\sub{e})^3 c}
\left(N\sub{e} \over \nu\right)^2 \left(m\sub{e}\over k T\sub{e}\right)^{3/2}
{\rm ln}\Lambda
\end{equation}
where it is assumed that the gas is singly ionized and $\Lambda$ is 
given by
\begin{equation}
\Lambda=\left(2k T\sub{e} \over \delta m\sub{e}\right)^{3/2} 
{4\pi \epsilon_0 m \over \pi\delta e^2 \nu} \approx
4.9573\,10^7 \left(T\over {\rm K}\right)^{3/2} {{\rm Hz} \over \nu}
\end{equation}
for $T\sub{e}<3.2\,10^5$K. The quantities $e$ and $m\sub{e}$ denote the
electron charge and mass and $c$ is the vacuum light velocity.

For a thermal plasma, the emission coefficient follows
from the Planck function
\begin{equation}
\epsilon(\nu)=\kappa(\nu)\times B_\nu(T\sub{e})\;.
\end{equation}

In the radiative transfer computations within the H~II region the
small frequency dependence of these continuum coefficients within the
molecular rotational lines is neglected.


\subsection{The Sobolev approximation}

As an initial guess to the level populations, the Sobolev approximation
may be used. This approximation is strictly valid for clouds with large 
velocity 
gradients  and velocities which are increasing outwards so that
all radiative coupling is restricted to a small spatial region. 
Then the level populations can be computed from local quantities only. 

The radiative energy density at a radius $r$ is determined by the
local source function $S_{j,l}(r)=\epsilon_{j,l}(r)/\kappa_{j,l}(r)$ 
\begin{equation}
u_{j,l}(r) = \left(1-\beta_{j,l}(r)\right)S_{j,l}(r) + \beta_{j,l}(r) I\sub{bg}(\nu_0)
\label{sobolev}
\end{equation}
where $\beta_{j,l}(r)$ denotes the probability for the escape of a photon
within the line $j$ from the cloud when emitted at $r$.
It can be computed from an integral over all spatial directions
from a given point by
\begin{equation}
\beta_{j,l}(r) = {1\over 2}\int_{-1}^1 {1-\exp(-\kappa_{j,l}(r)/Q(r,\mu))
\over \kappa_{j,l}(r)/Q(r,\mu)} \;d\mu
\end{equation}
with the line integrated effective absorption coefficient $\kappa_{j,l}$ and 
the velocity gradient in spherical symmetry
\begin{equation}
Q(r,\mu)=\mu^2{dv_r\over dr}+(1-\mu^2){v_r\over r}
\label{velgradient}
\end{equation}
Hence, the radial velocity gradient is the most
important quantity for the local energy density.

In contrast to the iterative solution of Eq.\ (\ref{balance}) it is possible
to apply a more efficient way here, since the energy densities are
analytically coupled to the level populations via Eq.\ (\ref{kappa}) and
(\ref{sobolev}-\ref{velgradient}). Using the derivatives of the matrix 
coefficients, the balance equations can be solved by a Newton-Raphson
approach. The corrections to the level populations within each step
will be computed from the matrix equation
\begin{equation}
\sum_i\left(\sum_k {\partial A_{jk} \over \partial n_i}n_k+A_{ji}\right)
\Delta n_i = \sum_i A_{ji} n_i
\label{diffbalance}
\end{equation}
where the $A_{ij}$ are the coefficients of the matrix of balance equations
(Eq. \ref{balance}).

\subsection{Computation of beam temperatures}

When the level populations are known, the beam temperature relative 
to the background is computed from the emergent intensity
by integration over the projection of the telescope beam on the
cloud.
\begin{equation}
T\sub{beam}(\nu)={c^2 \over 2 k \nu_0^2} \;{\displaystyle \int_0^{2\pi} d\phi
\int_0^\infty dp \;p\times (I\sub{surf}(\nu,p)-I\sub{bg}(\nu) )
f\sub{beam}(p,\phi) \over \displaystyle
\int_0^{2\pi} d\phi \int_0^\infty dp\;
p\times f\sub{beam}(p,\phi)}
\label{beam}
\end{equation}
The emergent intensity is the value on the cloud
surface $I\sub{surf}(\nu,p)=I_z(\nu,p,\sqrt{R\sub{cloud}^2-p^2})$.
A Gaussian profile for the beam is assumed
\begin{equation}
f\sub{beam}(p,\phi)= 
\exp{}\left(-(p-p\sub{offset})^2(1+\phi^2)^2 \over \sigma\sub{beam}^2\right)
\end{equation}
The projected beam width is computed from the angular
width by $\sigma\sub{beam} = \pi D^2/648000\; \sigma\sub{beam}['']$
where $D$ is the distance of the cloud. The standard deviation
$\sigma\sub{beam}$ is coupled to the full width of half power
by ${\rm FWHM}=2\sqrt{{\rm ln}2}\,\sigma\sub{beam}$. The program
computes a full map, i.e. the line profiles
at a given number of positions on a linear radial scan through
the cloud.

Within the program, all frequencies are treated in units of
$(\nu-\nu_0)/\nu_0$, velocities in units of $c$, and 
intensities in units of $\nu_0^3/c^2$, i.e. by a factor
$4\pi/h$ larger than the SI unit.

\section{The general code design}

The general design of the program is directed towards a high
accuracy of the computed line profiles. All of the errors in
the different steps of the program are explicitely user 
controlled by setting thresholds. All discretizations
necessary to treat the problem numerically are performed in
an adaptive way, i.e. there is no predefined grid at all
and all grid parameters may change during the iteration procedure.

Furthermore, the code was pushed towards a high flexibility,
i.e. the ability to treat a very broad range of all physical
parameters with the same accuracy and without numerical
limitations. E.g., the systematic velocities may range from
0 to several times the turbulent velocity depending on the
compiled frequency field size.
There is no inherent restriction to a distinct range of optical
depths. I have tested it for depths at line centre between about -5,
i.e. moderate masering and about 5000. However, the convergence
speed depends dramatically on the optical depths, the
level structure and the local velocity gradients.

The program is not optimized towards a high velocity of the
code. Although I got a factor of about 60 in velocity relative
to the precursor program by E. Kr\"ugel, other codes with lower inherent
accuracy may run another factor 10 faster. Nevertheless,
the code is suitable for an interactive work even on  a small
PC.

\section{User interface}

\subsection{Programm invocation}

The radiative transfer code is started from the command line by:

{\tt simline [}options{\tt ]}\\
where options can be:
\begin{description}
\item{}{\tt -fileinput}
\item{}{\tt -numdefaults}
\item{}{\tt -populateonly}
\item{}{\tt -stoponerror}
\item{}{\tt -overwrite}
\item{}{\tt -autonames}
\item{}{\tt -fits}, {\tt -block} or {\tt -lines}
\item{}{\tt -nocomments}, {\tt -unixcomments} or {\tt -vmscomments}
\item{}{\tt -writetau}
\end{description}

When the option {\tt -fileinput} is given, the program assumes that
all physical, observational and numerical parameters will be read
from files; interactive parameter input is switched off.
The option {\tt -numdefaults} sticks the set of the used numerical
parameters to the default values. With the option {\tt -populateonly}
the ray tracing part is skipped and only the level populations are
computed. The option {\tt -stoponerror} is intended for non-interactive
usage. By default the program asks for a correction of the input if
it detects some problem. With the option the code terminates in
case of errors.  The other options set the 
behaviour for the data file output (see Sect. \ref{datastruct}). 

All the options are intended to avoid additional interactive questions
during the program execution. If options are omitted, the corresponding 
parameters will be determined interactively during the program execution
except for the {\tt -writetau} option which is required if additional
optical depth profiles are to be generated.

\subsection{Terminal interaction}

The program is designed for an interactive computation of line
profiles so that the  flow is controlled by terminal
input. However, all parameters necessary for the computation
may be given either interactively or may be read from data files
except for the molecular constants which must be provided in a file.

There is a strict separation between parameters describing the 
molecule, parameters for the cloud geometry and physics, parameters 
of the observation, and numerical
parameters which control the accuracy of the different parts
of the computation.

A parameter input which was done interactively may be stored 
immediately in a file with the appropriate structure. This file can be read
at the next time of use, and modifications of single values
can be done by means of a text editor. In case of errors made
when interactively defining the data, the input procedure can be 
repeated by typing a character which cannot be converted to a numerical 
value at the next input question.

The two parts of the program may be run in arbitrary order
and repeatedly. At first one can compute either a new
cloud model, i.e. the level populations at all radial
points throughout the cloud, or one can read a file which
contains the results of a previous computation. Then, it is
possible to compute the observed line profiles for an arbitrary
number of lines or telescope parameters. For a hand-made line
fit one can change some cloud parameters and repeat the whole
procedure.

I recommend to use a graphics software in a parallel session/window
to get an immediate impression from the changes on the line profiles. 

\subsection{Data files}
\label{datastruct}

\subsubsection{Input files}

All input files are plain ASCII files so that they can modified with 
any text editor. Output is either plain ASCII or FITS to enable
easy examination, modification and use for further reduction in other programs.

The input files providing the molecular, physical, observational and 
numerical parameters consist of sets of dual lines. The first line in each
pair is treated as a comment line by the program and, in general,
contains a description of the parameter. The second line
contains the parameter in ASCII form. Exceptions occur in the file
for the molecular data where the levels, transitions and collisional
rates are given in tables with only one comment line per table
and in the cloud parameter file if the functions for the different
quantities are given by a data table which has also only one comment line.
I highly recommend to produce a first set of input files
using the SimLine program itself by interactive input of the
parameters. This guarantees the consistency of the produced parameter 
files with the structure expected when reading the parameters.
Starting from this first set it may be more feasible to change
distinct values in these files using a text editor. 

\subsubsection{Molecular data files}

Molecular data are read from ASCII files with a simple table structure.
Data files for several molecules accompany the SimLine executables and should
be used as general examples. 

The file starts with a comment line describing the molecule and ends with
an arbitrary number of comment lines for references. The entries
within the file are given as single values or tables each 
preceeded by a comment line explaining it. The following entries occur
\begin{description}
\item{}\tt Relative molecule mass:\\ \rm
 relative mass of the isotope given in atomic mass units
\item{}\tt Number of levels:\\ \rm
At most 50 levels can be read.
\item{}\tt Number of transitions:\\ \rm
At most 200 radiative transitions can be read.
\item{}\tt Levels: index, statistical weight, energy [cm\^{}-1], name\\ \rm
Here, a table with one line for each level has to be provided.
Index numbering should start by 1. To denote each level a string of 
8 characters may be used.
\item{}\tt Transitions: from index, to index, Einstein-As [s\^{}-1]\\ \rm
Here, a table with one line for each radiative transition has to be provided.
The indices have to correspond to the values in the level table.
\item{}\tt Collision rates\\
Number of temperatures\\ \rm
To separate the collision rates from the molecular radiation data
two comment lines are introduced here. Collision rate coefficients
can be tabulted for at most 10 different temperatures.
\item{}\tt Table of temperature values\\ \rm
Here, the temperatures have to be given - one per line.
\item{}\tt Number of transitions\\ \rm
At most 1000 collisional transitions can be read.
\item{}\tt Transitions from - to index,  rate coefficient\\ \rm
Here, a table with one line for each transition has to be provided.
In each line the indices of the two levels and the rate coefficients
for all temperatures have to be given.
\item{}\tt Comments and references\\ \rm
At the end of the file an aritrary number of comment lines
can be used where references to all data should be provided.
\end{description}

Alternative to the native files for molecular data, SimLine can read
LAMDA files in the RADEX format. However, it is currently limited to 
one collision partner so that LAMDA files with more partners are
rejected to avoid ambiguities.

\subsubsection{Output files}

Two types of result files are produced. The first one contains the
level populations calculated as the result of the lambda iteration.
This file is only intended for further processing by SimLine in the 
computation of line profiles.
There are no comments in the file. The first line contains the
file name of the molecular data file. In the second line we find
the number of radial points, the number of
rotational levels treated and the background radiation temperature
at the location of the cloud. In the following block, the physical
parameters at the different radial points are stored.
Each line contains the number of the radial point, its radius,
the molecule density, the systematic velocity, the 
thermal+turbulent velocity $\sigma$, and the correlation length over
$\sigma$. A second block contains the lines with the level 
populations. For each 
radial point (denoted by its number) the populations at all 
computed levels are stored.
The same file structure as produced here 
must be given when reading a precomputed cloud model from a file in order 
to calculate new line profiles.  

The second file type contains the line profiles. These files are designed
to be readable by graphics programs as IDL, GREG or GNUPLOT. Two general 
file structures are possible - FITS and ASCII. In case of ASCII files
one can further select whether to write the data points as a compact
table or each point with offset and velocity in a separate line.
In FITS format the lines are stored in a two-dimensional array
with the first axis given by the angular offest and the
second axis by the velocity. In ASCII block mode
the two dimensional field of beam temperatures
is stored as a block where the lines represent the different beam
offsets and the columns the frequency/velocity points. In two preceding
comment lines the offset and frequency values are given.
The line-oriented structure requires much more disk space but can be read by 
more programs. Here, each line contains three numbers - the beam offset,
the frequency, and the beam temperature. Consequently the file
contains (number of offsets)$\times$(number of frequency points) lines.
Preceeding comment lines specify these numbers.
The format of comment lines can be chosen to be in VMS or UNIX format,
i.e. whether a `!' or `\#' shall be used to indicate the beginning of 
a comment line. Furthermore one can also skip the comment lines.

In addition to the line profiles, profiles of optical depth can
be written by providing the {\tt -writetau} option. The optical
depth files have exactly the same structure as the line intensity files.
With the option one optical depth profile is written for every line.

The default parameters for the structure of the line profile files
and the possibility to always overwrite existing files with newer
ones can be specified as options on the command line.
The options {\tt -fits}, {\tt -block} and {\tt -lines} turn on FITS
format, ASCII block mode or ASCII line 
mode for the file structure. The options {\tt -nocomments}, 
{\tt -unixcomments}, and {\tt -vmscomments} define the format of the 
comment lines in case of ASCII files.
When the additional option {\tt -overwrite} is given, existing files 
will be replaced without further notice. Otherwise, every overwrite 
has to be confirmed interactively.
With the option {\tt -autonames} the program does not ask for names of
output files but creates those names from the name of the input file
appending a string for the considered transition, based on the level names.
All these options can be given also on compile time via conditional
defines (e.g. {\tt -Dnocomments -Doverwrite}). 
If no parameters are specified, they will be asked interactively before
each file output. 

The units in the files are always pc for radii and km/s for
velocities/frequencies. Beam offsets are given in degrees on FITS 
output and in arcseconds in case of ASCII output. 
Frequencies above a line central frequency, i.e. blue
shifted parts are counted as negative velocities and red shifted parts
positive.

\subsection{The input parameters}

\subsubsection{Physical cloud parameters}

\begin{description}
\item{}\tt Background radiation temperature [K]:\\ \rm
For normal galactic sources, the cosmic background with
a temperature of 2.73 K is dominating the external radiation
in the wavelength range of the rotational transitions. Looking
at early universe objects one should adapt higher values.
\item{}\tt Number of shells for the cloud representation:\\ \rm
The cloud may be subdivided into a number of shells with different
laws for the physical parameters. These shells are only
used to to define the input data and have no relation to the
radial shells later used in the radiative transfer computation.
It is possible to specify discontinuities at shell edges 
in case of power law regions and file input, however, they
should be reserved to the experienced user. The maximum number
of shells is restricted to about 180.
\item{}\tt Use shells as power law regions or as data table ? [0/1]:\\ \rm
Within the shells one can either specify the radial dependence of
the physical parameters by their power law exponents or in terms 
of a data table where the corresponding shell behaviour is interpolated.
\item{}\tt Radius of the inner core [pc]:\\ \rm
To avoid the singularity of power laws at zero, the user has
to define an inner core. Within this core all parameters are assumed
to be constant. Furthermore, the core may consist of an H~II region.
In case of a data table input the core radius is identical to the
innermost table radius.
\item{}\tt Electron density [e/cm\^{}3] and temperature:\\ \rm
In case of an H~II region for the inner core, these parameters
have to be specified to compute the continuum absorptivity and
emissivity. An electron density 0 means that no central HII region 
will be assumed.

\end{description}

\noindent{}For each shell the following parameters have to be specified.
\begin{description}
\item{}\tt Outer radius of the shell [pc]:
\item{}\tt Hydrogen density [H2/cm\^{}3]:
\item{\tt} Temperature [K]:\\ \rm
The kinetic temperature determines the collision rates and provides
a contribution to the local line width.
\item{}\tt Relative molecular abundance [X/H2]: \\ \rm
Here, possible molecule depletions should be taken into account.
\item{}\tt FWHM of turbulent velocity [km/s]:\\ \rm
The thermal velocities will be added to get the total kinetic line width.
\item{}\tt Turbulence correlation length [pc]:\\ \rm
The scale on which the autocorrelation function of turbulence
goes to 1/e.
\item{}\tt Systematic radial velocity [km/s]\\ \rm
Infall is counted with negative numbers.
\end{description}
In case of power law input, all parameters, except the radius, are 
characterized by the value at the inner boundary of the shell and a 
radial exponent. In case of data table input only the values at the
corresponding radial points have to be given.

\subsubsection{Observational parameters}

\begin{description}
\item{}\tt Upper level of the transition to be observed:\\ \rm
By default, only one line is considered at a time. Here, we need the number of
the upper level in the molecular data file.
\item{}\tt Lower level of the transition to be observed:\\ \rm
Number of the lower level in the molecular data file.
If zero is given both for the upper and the lower level, SimLine will
compute the line maps for all excited transitions of the molecule 
using the same observational parameters.
\item{}\tt Frequency resolution [km/s]:\\ \rm
This might be restricted by the detector or the purpose of the observation. 
\item{}\tt Beam width [FWHP in ``]: \\ \rm
The beam is considered to be Gaussian.
\item{}\tt Step size for mapping [``]: \\ \rm
One typically choses steps in the order of half a beam width.
\item{}\tt Central offset for the first point [``]: \\ \rm
In case of observational offsets or composed maps, one may start 
the radial map not at the centre of the cloud but with a finite offset.
\item{}\tt Number of map points [0=map whole cloud]: \\ \rm
A number 1 means just one observation, not a map. When giving a
zero, the program automatically computes the number of map points
up to the outer radius of the cloud.
\item{}\tt Distance of the cloud [pc]: \\ \rm
Needed to translate the physical cloud size into angular units.
\end{description}

\subsubsection{Numerical parameters}

The same set of numerical parameters can be used to determine the 
accuracy in the computation of the level populations and in the computation
of the emergent line profiles. In a few cases a parameter is interpreted
slightly different in the two parts.

\begin{description}
\item{}\tt Cut-off for Gaussians in the line emission [sigma]:\\ \rm
This one and the next parameter determine at which argument
a Gaussian is considered to be zero. For higher arguments
the analytic treatment is truncated. The cut-off in line emission
determines at which frequencies in radiative transfer the absorptivity/
emissivity is treated as zero, i.e. the background field remains unchanged.
The default value implemented in the program is 3.2. 
\item{}\tt Cut-off for Gaussians in freq.\ integration [sigma]:\\ \rm
This parameter says where to set the limits in integration over
the intensity multiplied with a Gaussian profile. It is applied in
the frequency integration to determine the local energy density at each 
point and in the spatial integration over the telescope beam.
The default value is 4.2.
\item{}\tt Resolution in scanning the Gaussians [sigma]:\\ \rm
The parameter defines the resolution of the grid for scanning the
frequency scale in the integration of the radiative energy density. 
This frequency grid is locally adapted to be always
at least as fine as this resolution factor times the local $\sigma$.
The default value is 0.55.
\item{}\tt Maximum change factor for the level densities between grid points:
\\ \rm
The gas densities and level populations are linearily interpolated 
between the radial grid points. This parameter determines the maximum 
change of the level densities $n_j$ between the grid points. In case of
steeper changes, additional grid points are added. The factor has to be 
greater than 1. The default value is 1.3.

In many cases, the reduction of this parameter provides the strongest
improvement of the computational accuracy at the costs of a strong
increase in the time consumption. It may, therefore, be considered as 
the most important numerical parameter.
\item{}\tt Hysteresis in the radial grid adjustment [0..1]:\\ \rm
Since the radial grid is dynamically adjusted to the gradient of the
level densities $n_j$, grid points may be added or removed during the
lambda iteration. The hysteresis gives the relative difference between
the thresholds for addition and removal of points. It
supresses grid oscillations during the iteration.
The default value is 0.2.
\item{}\tt Maximum change of the radial velocity between grid points [sigma]:
\\ \rm
Beside the density structure, also the velocity structure may require
additional radial grid points. The systematic velocity is linearily interpolated
between the grid points. The parameter determines the maximum change
of the systematic velocity between two radial grid points. In case
of stronger changes, additional grid points are included.
The default value is 2.5.
\item{}\tt Maximum relative distance of points on the spatial integration scale:
\\ \rm
This parameter determines the number of displacement values treated
in the radiative transfer code. Since the computation of the energy density
at a given radial point requires the integration over the $z$ axis (Eq.\ 
\ref{uzint}), many points should be given along the $z$ axis. The parameter
constrains the maximum relative distance between two points on the $z$ axis
within this integration procedure.
The default value is 0.25.
\item{}\tt Maximum velocity shift between radiative transfer points [sigma]:
\\ \rm
The radiative transfer code interpolates the emissivity and absorptivity
for each frequency linearily within each step. In case of large velocity 
gradients, this will be inappropriate since the line profiles are
widely shifted between two grid points. Here, additional steps have to
be included. The parameter gives the maximum length of such a step on
the scale of velocity in $z$ direction $v_z$. It is certainly a good idea
to take the same value as for the frequency resolution so that the
frequency shift between consequent radiative transfer points is always
according to at most one grid spacing at the frequency scale.
The default value is 0.55.
\item{}\tt Initial populations [1=thermal,2=no radiation,3=Sobolev,
  4:file input]:\\ \rm
Here, one can determine the first guess for the level populations the
lambda iteration starts with. 
The thermal distribution is an appropriate approach for clouds
which are optically thick within at least one line of the molecule
considered. The radiation free distribution should be chosen for
optically thin clouds. The Sobolev approximation is good for clouds
with large velocity gradients. However, I found that, due to the high
efficiency of the accelerated lambda iteration method, the influence
of the initial guess on the total number of iterations is relatively small.
The level populations output file from one model run can be used
as an input file for another model. This is especially useful when
iteratively fitting observations by scanning the parameter space.
The default value is 2.
% Undocumented feature: negative values skip the whole iteration.
\item{}\tt Neglection threshold for the level populations:\\ \rm
The parameter sets the truncation of the equation system for
the level populations. Whenever the excitation of one and all higher
levels troughout the whole cloud falls below this threshold, the
level is neglected in the further computation. According to the
machine precision used for the representation of the numbers
values below $10^{-14}$ make no sense.
The default value is $10^{-8}$.
\item{}\tt Relative accuracy as convergence criterion:\\ \rm
This is the convergence criterion. It is tested only after a full 
lambda acceleration cycle. When even the lambda acceleration does
not find a larger step size than this accuracy value, convergence
is reached. To obtain accuracies better than $2\,10^{-7}$, the
program has to be compiled with {\rm real*8} numbers instead of
the normally used {\rm real*4}. The default value is $10^{-6}$.
\end{description}

Due to the strong nonlinearity of the whole line radiative transfer
problem it is not possible to give a general error estimate for
a given set of numerical parameters. Tests using a broad range of
cloud parameters have shown, however, that the error in line intensities
using the default numerical parameters falls always below 5--6\%.

When the command line option {\tt -numdefaults} is given at program 
invocation, the default numerical parameters are used in all computations.
This can avoid the annoying confirmation of the usage of these default
values before every computation in cases when one does not plan to change
them. Then, it is, however, impossible to use changed numerical
parameters within the same program run when it turns out that a
given problem is not soluble with these default numerical parameters.

\subsection{Error messages and warnings}
\label{warnings}

Due to the static declaration of fields in FORTRAN 77 the possible
numbers of grid points on each scale is restricted to the maximum
size of the corresponding fields. There are two different precompiled
binary versions available. In the normal version the field size for the
radial points is set to 180, for the impact parameters to 270,
for the frequency points to 223, for the molecular levels to
20 and for the number of radiative transitions to 19. This version
should be suitable for most linear molecules. Most errors occur
when the combination of cloud properties and numerical parameters
requires a grid which exceeds these numbers in one dimension.

If the radiative transfer problem requires some
field size larger than the compilation values you will face an error 
message and either a return to the main input loop (default)
or a stop of the program with the exit code 1 (if the
option --stoponerror is set in the program call).
You may find the following messages:

\begin{description}
\item{}{\tt Number of radial points insufficient to treat the problem
with the required\\ accuracy !} \\
The gradient in the level densities requires more than
140 radial points either due to a very steep density gradient or due
to strong changes in the level populations, e.g. at the edge
of very thick clouds. The user will be brought back to the
input routines to change some parameters. One might either 
somewhat increase the numerical parameter characterizing the
maximum change of level densities between grid points or in case
of extremly differing densities within the cloud substitute
regions with steep density gradients by a discontinuity in the gas
density.

\item{}{\tt Number of ray points insufficient to treat the problem
with the required accuracy !} \\
This error may scarely happen when
the user has chosen a very small distance of the grid points on
the $z$ integration scale in connection with many radial grid points.
It may also appear in the computation of the emergent line profiles
when the systematic velocities are very large compared with the local
random velocities or a large cloud is mapped with a very fine beam.
He will be brought back to the input routines. One should slightly
increase the numerical parameter characterizing the relative distance
of grid points on the spatial integration scale or, in case of
the large velocities, increase the numerical parameter giving the
resolution in scanning the Gaussians.

\item{}{\tt Number of frequency points insufficient to treat the problem
with the required\\ accuracy !}\\
Looking at the maximum velocities in 
the cloud and the frequency resolution required either by the local line 
width in some points of the cloud (first part of the program) or by the 
observation (second part), more than 223 frequency points would be required.
The user will be brought back to the input routines to change some 
parameters. When the error appears in the first part, i.e. the computation
of the level densities, one should increase the numerical parameter
characterizing the resolution in scanning the profiles or slightly
reduce the numerical parameter for the cut-off in the frequency 
integration. However, there is no solution when the systematic
velocities exceed the local line width by many order of magnitude.
If the error appears in the second part, i.e. the computation
of the emergent line profiles, one should consider to reduce either the
frequency resolution chosen as observational parameter or to reduce 
the numerical parameter for the cut-off in the frequency integration.

\item{}{\tt Numerical singularity occured for the given parameters !}\\
This message occurs when the general ``worst case check'' fails indicating
that the code is no longer able to compute non-singular values for
the level populations. This error typically results from unphysical
input parameters like negative cloud temperatures or temperatures
high above the range of valid collision parameters or negative radii.
One should carefully inspect the input parameters to detect the 
true reason.
\end{description}

Other error messages may appear in case of problems with the file
input/output.
\begin{description}
\item{}{\tt  I cannot open the input file}\\
This message is thrown if the specified input file does not exist or is
not readable. With the option --stoponerror, the program stops with
exit code 2.
\item{}{\tt Format error in input file}\\
This message is thrown if the line structure of the input file
cannot be recognized, i.e. the order of parameters is wrong or
comment lines are missing. With the option --stoponerror, the program stops with
exit code 3.
\item{}{\tt I cannot open the output file}\\
This message is thrown if the output file exists and the
--overwrite option is not set, the output file exists and
the permissions do not allow to overwrite it or the output
file cannot be written because a non-existing directory is specified.
With the option --stoponerror, the program stops with
exit code 4.
\end{description}

In the input routine for the molecular data file it is checked whether
the parameters in the file fit into the internal input fields
regarding the number of levels, radiative transitions, temperatures,
and collisional rate coefficients. However, if the file contains
more levels or transitions than the radiative transfer code can treat, 
the molecular system will
be simply truncated at high energies by virtually removing high
levels until all transitions and levels fit into the fields.

There is no explicite test whether the maximum number of atomic 
levels treated by the program is sufficient for a given problem.
I case of doubt one should inspect the output file for the level
densities looking at the excitations of the last level. The number of 
levels can be a serious restriction when treating relatively warm
and optically very thick clouds.

A small number of consistency checks is done for the parameter input
and the following messages may appear:

\begin{description}
\item{}{\tt Line overlap detected: transition - from - to - at - Hz}\\
At the moment, the program only checks for overlapping lines
but does not yet treat them in radiative transfer. It is just a warning
that the results will not be reliable in this case.

\item{}{\tt There is no excited transition from level - to level - !}\\
This error occurs when the observational parameters select a transition
that does either not exist as radiative transition in the molecular 
data file or that includes a level
the excitation of which falls below the neglection threshold of
the previous computation of the level densities.
A new transition has to be selected.

\item{}{\tt Only finite beam widths are possible !}\\ The FWHM of the beam 
has to be $>0$.
\end{description}

Nevertheless, most parameters are not checked and strange results
or abnormal program termination may occur in case of unphysical
or numerically absurd parameters. E.g. one has to take care in the
input parameters that the kinetic temperature in all parts of the cloud
should be above the cosmic background temperature. Otherwise, the
code will probably abort computation at some point.

\section{Hints for the application of the code in line fitting}

My own tests and the experiences of some colleagues lead to a few recommentations the user of the program should be aware of:

\begin{itemize}
\item{}To get a first feeling one should always start using the
default numerical parameters directly implemented in the program.
Only in a late stage, the accuracy of the result should be checked
by setting stronger limits to the numerical accuracy.
\item{}When fitting a number of lines, the user should always start
with a simple one-shell-model (except there is a detailed physical model
in background suggesting a more complicated structure). Modifying
the parameters within this one shell will already provide the possibility
to produce a very wide variety of line shapes and intensity ratios.
In complicated models, one can easily get lost within the
parameter space. {\it Remember, in 1-D, there is never only one unique
solution to a line fitting problem !}
\item{}To get a reasonable initial guess for the  cloud parameters,
traditional ways for the determination of the gas temperature and the
column density, like the comparison of the line strengths from different
isotopes of a considered molecule, should always serve as starting point.
\item{}Whenever the cloud model covers a range of gas densities $>10^3$
or the density structure is relatively complex, it may happen
that one runs out of the number of points available for the radial grid. 
For smooth density structures, the user should relax the numerical 
parameter giving the maximum difference in the level densities between 
two radial points. In case of strongly variing density profiles it 
may be more appropriate to substitute steep gradients by artificial 
jumps in the density (The input of discontinuities in the parameters 
is only possible using file input).
\item{}At the end of a density fit, the user should check the
consistency of the temperature estimate with the computed line temperature
of the thick line used for the temperature derivation. I found that
even quite thick lines ($\tau > 10$) are often not completely 
thermalized. At least when ${}^{12}$CO is used, this way can provide a much 
more accurate temperature determination.
\item{} As a first orientation to constrain the turbulence parameter
space, one might follow Miesch \& Bally (1994) which got correlation lengths
of about 0.1 pc for the best resolved clouds. Their values for the Kolmogorov
exponent $\gamma$ giving the radial dependence of the turbulent
velocity dispersion range from 0.74 for L1228 down to values around 0.5
for relatively many objects. For a homogeneous isotropic infinite 
turbulence, there is also a theoretical relation between $\gamma$ and 
the correlation length: $\gamma \propto 1/\ln(r\sub{c})$ showing a general
anticorrelation. 
\item{}The approximation of effective absorption coefficients for a 
clumpy medium is valid both for clumping in space and velocity space.
The used description by the correlation length is according to a
homogeneous medium with pure turbulence in velocity space. There is
no additional parameter for the filling factor in space
in case of a clumpy medium, since this parameter is not independent.
Instead, one has to use the clump hydrogen density for the hydrogen 
density parameter,the molecular abundance has to be reduced by the 
filling factor and the correlation length has to be increased by one 
over the filling factor.
\end{itemize}

\section{System requirements}

The program is written very close to standard FORTRAN 77 so that most 
f77 compilers should be able to compile it.
It has been tested as it is on DEC Alpha running OSF/1, DEC VAX
running Ultrix, SUN Sparc running SunOS, IBM RISC running AIX, HP
with HP-UX, and Intel machines running Linux (with f77(f2c), g77, 
and Absoft-f77). 
On DEC VAX machines running VMS, four lines in {\tt ltrio.f} have to be 
changed due to a different record treatment of this operating
system compared to the Unixes above. On none of these platforms
the compiler or linker should give any warnings. I did not yet succeed to
run the code under MS-DOS due to a strange memory management of the
Microsoft-Fortran-Compiler/Linker.

The memory usage is mainly determined by the one large field
of energy densities at all radial points. In the normal precompiled 
version it consists of $180\times 540\times 19$ real*4
numbers which makes in total about 7.4MB alone. In the extended
version it uses $220\times 660\times 90$ real*4 numbers, i.e.
52 MB.  Adding another 10\% for all other fields
and the dynamically linked libraries gives a reasonable guess for 
the total memory consumption.

\section{Limitations and bugs}

One should keep in mind that a main limitation in all radiative 
transfer computations comes from the accuracy of available
collision rates and the temperature range where they are given.
At temperatures outside of the interval given in the molecular data
file the collision coefficients will be linearly extrapolated.  
Moreover, one should also consider that the maximum number of levels
is restricted in the code so that high temperatures combined with
high densities might also
provide a non-vanishing excitation at levels which are neglected
in the computation. The user has to check for his own, that the 
maximum level number is sufficient.

The program is not designed to fit the situation of a strongly variing 
turbulence. The gradient in both the absolute value of the turbulent 
velocity dispersion and in the turbulence correlation length should 
not exceed the steepest gradient of either 
the systematic velocity or the density. 

Due to the limited number of frequency points the systematic velocities 
may fall at most one order of magnitude above the random velocities.

At the moment, I know of no bugs in the code.

\section{Future plans}

New versions of the program are planned to contain:
\begin{itemize}
\item{} interaction with dust continuum radiation
\item{} collisional excitation by multiple species
\item{} improvement of the treatment of scale-invariant turbulence 
\item{} correct treatment of overlapping lines
\item{} massive parallelization
\end{itemize}

\section{Internals - Structure of the code}

Here, I can only give a short overview over the general structure of
the code. It shall help as a guide through the source code in case
of a deeper inspection of the numerical internals or in case of
error debugging. The overview explains the function of the different
routines and how they are concatenated within the program. It cannot
substitute for a study of the sources which I have tried to comment
relatively detailed.

Please, do not perform any essential changes to the sources without 
contacting me and send me all your suggestions to improve the code or
to make it more user-friendly.

\subsection{Main program}

LINETRANS is the main program. It provides the loops for calling
MICRO and LINE, the routines for the computation of the molecular
level populations and the observed line profiles.

Both MICRO and LINE start with reading the parameters specified
by the user. MICRO then initializes the grids and starts the
loop made of the radiative transfer (STPMIC), the lambda iteration
(LAMBDA, LEVMIC) and the grid adjustment (ADAPTGRID). In case
of convergence the results may be stored (WRITEDENS).

LINE continues with setting up the grids for beam offset and frequency,
calls the radiative transfer in one line (STPLINE), the beam 
integration (INTBEAM) and finally writes the results (WRITELINE).

PRESENT shows the opening and closing screen of the program.

\subsection{Input - Output}

The routines for the parameter input READPHYS, READNUM, READOBS
have the same general structure. One part runs the file input
and another one the interactive input. In case of interactive
input an inverse version of the file read can be used to produce
a parameter file.
For the input of numerical parameters (READNUM), there is an
additional part so that it is
possible to use the built-in default values or to use the 
parameters from a previous computation if called more than once.

The routine for the input of the molecular data file READSPECIES
allows no direct terminal input and contains a translation from
input fields which should be large enough to read all reasonable
molecular data files and the internal fields which are used
in radiative transfer.

WRITEDENS is the routine to store the level populations in a file 
with the structure described in Sect.\ 2.2. READDENS reads the
populations from such a file. Finally, WRITELINE is used
to store the integrated beam temperatures in a file with
variable structure. In case of FITS format for line output
it calls SAVEFITS which needs to be linked with the FITSIO
library and the auxiliary routines PRINTERROR and DELETEFILE
from fitsadd.f. The default parameters for the structure
of the files written are set in PREDEFINE which reads the command 
line. The function IOERROR is called in case of
file-I/O exceptions. It produces the error messages and determines
the change of the program flow.

\subsection{Radiative transfer for all lines}

These routines perform the radiative transfer computation
necessary to obtain the radiation energy density at each
point and within each molecular transition.

STPMIC contains the loops for all displacement parameters
and all radial points. For each ``ray'' LEFTKPL is called
to initialize the fields for the intensity (with the background
radiation - UBACK) and the absorption coefficient 
at the edge of the cloud.

To carry out the transfer between two radial grid
points, at first FIRSTSTEP is called. This subroutine determines
the number of real radiation transfer steps necessary between
these points to have sufficiently small shifts in the line profiles.
If more than one step is required NEXTSTEP is called in a following loop
to compute the size of further steps.

TRANSFER is the real radiative transfer routine. It constrains
the frequency range where the line profile is non-vanishing
and computes the new intensity in this range. In case of small
optical depths a linear approximation of the exponential functions
(Eq.\ \ref{transfer}) is done, otherwise the full function will
be treated assuming that $\kappa$ and $\epsilon$ show a linear
dependence within the integration range. CTRANSFER computes
the continuum line transfer within a central H~II region. It assumes,
that the opacity and source function remain constant within this core. 

At each radial point STPMIC calls the integration routines
UNYINT and UZINT to compute the energy density.

\subsection{Radiative transfer for a single line}

The subroutines STPLINE, FIRSTLSTP, NEXTLSTP, LEFTKPL, TRANSFL, and
CTRANSFL are reduced versions of STPMIC, FIRSTSTEP, NEXTSTEP, LEFTKP, 
TRANSFER, and CTRANSFER in the sense that they compute the radiative transfer
only in a single line instead of all transitions important for a 
molecule. They are used to compute the emergent line profiles
which are to be observed. No intensity integration is called.
The additional routine TAUCENTRAL is used here to give a control
output of the central opacity in the line centre. 

\subsection{Initialization and adjustment of the grids}

INITRAD builds the initial radial grid. It checks the gradients in density 
and velocity and takes as many points as necessary to fulfil the
conditions for the maximum changes specified by the numerical 
parameters. Additional points are included in case of discontinuities
detected in the density structure. To assign the physical parameters
density, temperature, systematic velocity \dots{} to a radial point
VALUES is called which computes these values from the physical input
parameters.

INITN initializes the level populations at the different radial grid
points according to the chosen first guess approximation. In case of
a central H~II region, it calls UFROMHII which computes the integrated
intensity at a given point due to the continuum radiation from
the H~II region.

FILEDENS takes both the initial radial grid and the level populations
from an procumputed file using READDENS. It thus substitutes both 
INITRAD and INITN in case of file input for the initial guess.

INITFREQ scans through all points and determines the local
line width. From this quantity the maximum frequency/velocity step
size for each point is computed. OVERLAP checks for line overlap
in the given velocity regime.

ADAPTGRID is responsible for changing the radial grid so that
the differences between the points fulfil the conditions for
the maximum change of the level population between two points and the
maximum velocity difference. It calls ADDPOINTS which checks for the need
of additional points and includes them and REMOVEPOINTS which
seeks for superfluid points which will be removed. In case of adding
new grid points VALUES is called which computes the gas density and
velocities at the new radii.

REMOVELEVELS checks whether some levels may be removed due to
negligible population and calls LINFROMLEV in that case to
reaarrange the field of transitional coefficients.

The routines MAKEPS and PSFORLINE build the grid of displacement 
parameters providing the ``rays'' on which the radiative transfer is 
calculated. MAKEPS is responsible for the grid used in STPMIC,
PSFORLINE builds the grid for STPLINE. Both grids differ.
MAKEPS controls that for a given radial point the relative distance
of the displacment points on the z scale is below the 
numerical parameter chosen in NUMREAD (see Sect.\ 2.3.3). 
The grid is made as fine as necessary to fulfil this condition.
PSFORLINE starts with the radial grid and makes it denser for
the displacement grid to have the same accuracy in covering the beam.
In contrast to MAKEPS all radial points are used.

\subsection{Solution of the balance equations}

MATRIXUP is the routine to compute the coefficients and set up
the transition matrix. Together with the normalization equation 
the transition matrix forms the linear system of 
balance equations (Eq.\ \ref{balance}).
LUSOLVER solves this system calling LUDCMP for a LU decomposition,
LUBKSB for the back substitution and MPROVE for the iterative 
improvement of the solution.

LEVMIC calls either ULEVMIC or DLEVMIC which store the
the energy densities or the level populations at all radial points 
for three cycles to be used in the lambda acceleration. After three 
ordinary lambda iteration cycles LAMBDA or DLAMBDA is called which 
compute a second order lambda acceleration step to the energy 
densities  or the level populations.
In case of a sufficiently small acceleration step 
(convergence criterion) LEVMIC terminates the iteration. HIISMTH
corrects the obtained level populations within the HII region
to enable a uniform treatment for all grid points.

KAPSRC is a small routine to compute the new absorption coefficients
and the source functions from the level populations using the
Einstein coefficients from AIJ,BIJ, and BJI. EFFKAP is the routine
which obtains the effective absorption coefficients for the turbulent
medium from the correlation length and the clump absorptivities.
ATAU is the corresponding reduction function. In case of a central
HII~region, HIIKAP corrects the opacities for the most inner point
to get the continuum values.

\subsection{The Sobolev approximation}

SOBOLEV contains the loop for all radial points and all iterations
to determine the local energy densities and level populations in
the Sobolev approximation. It computes the velocity gradients
which are used by ESCAPE to determine the escape probability 
in a given direction and its first derivative with respect to $n_j$. 
GAUSSOB integrates the escape probabilities over all directions
using a Gaussian quadrature, the grid points of which are 
initialized at the beginning by GAULEG. From the integrated
escape probabilities SOBOLEV computes the local energy densities
which are used by SOBMATRIX to obtain the level populations.

SOBMATRIX is the routine to set up the differential form of the
balance equation system used for the Newton-Raphson approach
to the local level populations (Eq. \ref{diffbalance}). It calls AMATRIX to get the
original form of the matrix (used on the right hand side of the system)
and DMATRIX for the differential form (for the left hand side).
The system is resolved by LUSOLVER. SOBKAPSRC is the routine
to get the absorptivities and the source function and their derivatives
from the level populations. SOBEFFKAP computes the effective absorption
coefficient and its derivatives for a clumpy turbulent medium.

\subsection{Intensity integration}

UNYINT sets the limits for the frequency integration at each
radial point and multiplies the intensities with the local
line profile (Eq.\ \ref{uzint}). Then, it calls QUINTGAUSS, 
a simple equidistant integration routine, which carries out the
frequency integration. The line profile is computed by PHI.

UZINT prepares the fields for the $z$ integration (Eq.\ \ref{uzint}).
In ranges with very small steps it calls the linear integration
routine DINTLIN, for larger steps it calls the cubic spline integration
routine DINTCUB. DINTCUB itself makes use of the cubic spline 
interpolation routine SPLINE.

INTBEAM is responsible for the intensity integration over
the telescope beam. It computes the range of the cloud covered
by a Gaussian beam and calls SPECERR for the $\phi$ integration
and finally DINTCUB, the cubic integration program, for the
$p$ integration. 
SPECERR computes an numerical approximation to the $\phi$
integral within Eq.\ (\ref{beam}) which is independent from
the intensity.

TRANSLAT is a small routine translating the angular beam width
of the telescope into the projected beam width on the cloud.

\subsection{Molecular constants}

The following routines contain the full ``physics'' of the problem.
All physical quantities are specified within them. All previous
routines only contain ``numerics''. The fields used here are
read in advance by READSPECIES.

INITEINSTEIN looks which molecule shall be considered and initializes 
the fields containing the molecule specific quantities, especially
$a_{j,j-1}$, $b_{j,j-1}$, and $\nu_0$. The functions AIJ, BIJ, BJI
and CFREQ simply read the tables of $A_{j,l}$, $B_{j,l}$, and $\nu_0$
and provide the results as function values.

PROBTHERM, VTHERM, UBACK also base on the tables produced by
INITEINSTEIN. These functions provide the population of a distinct 
level in case of thermal distribution, compute the mean thermal velocity
of the molecules at a given temperature, and determine the intensity
of the cosmic background radiation at the frequencies of the
molecule transitions. HIIINIT computes the continuum opacity and
source function for all transitions of the considered molecule
from the table of frequencies and the electron density and temperature 
within a given HII region.

CIJ obtains the collision rate coefficients at a given temperature 
for the transition of interest. It calls
ILOCAT to find the appropriate table entries and interpolates 
the values for the given temperature. 

\section{Source files}

The following files contain the sources for the different routines:

\begin{description}
\item{} {\tt ltr.f}: LINETRANS, MICRO, LINE, PHI, TRANSLAT, PRESENT
\item{} {\tt ltrio.f}: READPHYS, READNUM, READOBS, WRITEDENS, READDENS, WRITELINE,\\
SAVEFITS,IOERROR, PREDEFINE
\item{} {\tt initial.f}: INITRAD, INITN, FILEDENS, INITFREQ, 
ADAPTGRID, ADDPOINTS,\\ REMOVEPOINTS, VALUES, UFROMHII, OVERLAP
\item{} {\tt stpdiff.f}: STPMIC, FIRSTSTEP, NEXTSTEP, LEFTKP, TRANSFER, MAKEPS
\item{} {\tt stpline.f}: STPLINE, FIRSTLSTP, NEXTLSTP, LEFTKPL, TRANSFL, PSFORLINE,\\ TAUCENTRAL
\item{} {\tt matrix.f}: LEVMIC, ULEVMIC, DLEVMIC, LAMBDA, DLAMBDA, 
REMOVELEVELS,\\ LINFROMLEV, KAPSRC, EFFKAP, ATAU, HIIKAP, 
HIISMTH, MATRIXUP, LUSOLVER, LUDCMP, LUBKSB, MPROVE
\item{} {\tt sobolev.f}: SOBOLEV, ESCAPE, GAUSSOB, GAULEG, 
SOBMATRIX, SOBKAPSRC,\\ SOBEFFKAP
\item{} {\tt einstein.f}:  AIJ, BIJ, BJI, CFREQ,  
PROBTHERM, VTHERM, PLANCK, UBACK, HIIINIT
\item{} {\tt collrate.f}: READSPECIES, CIJ, ILOCAT
\item{} {\tt ltrinteg.f}: INTBEAM, SPECERR, UNYINT, QUINTGAUSS, UZINT, DINTCUB, 
DINTLIN,\\ SPLINE
\item{} {\tt fitsadd.f}: PRINTERROR, DELETEFILE
\end{description}

\section{Auxiliary programs}

I have written some small auxiliary programs, mainly for testing purposes
around the SimLine code. The program TESTCLOUD reads the level populations
and energy densities from a data file produced in the first part of SimLine 
and computes the excitation temperatures, radiation temperatures
and few additional quantities. The input file must have the structure
as produced by WRITEDENS, the output file is self explanatory.

The program CLINE reduces the two dimensional data set as produced by
WRITELINE and extracts the profile for zero beam offset, i.e. the
resulting file contains only one line profile. This may be suitable for
users which do not have a sophisticated graphics front end as GREG
or IDL.

\section{Acknowledgements}

A direct precursor of this program was written by Endrik Kr\"ugel 
(MPIfR Bonn). Basing on his code of about 1000 lines, I could set
up this more comprehensive and flexible program. The subroutines 
CIJCOH2 and ILOCAT are still very similar to the original code. 

The local statistical treatment of clumping or turbulence roughly 
follows the formalism derived by H.M. Martin, D.B. Sanders \& R.E. Hills
(1985).
The accelerated lambda method was published by Lawrence Auer
in 1987. The LU decomposition method, the Gaussian quadrature, and 
the spline interpolation algorithm were taken from ``Numerical Recipies''
(Press, Flannery, Teukolsky \& Vetterling 1986).

\section{References}

\begin{description}
\item{}Auer L.H. 1987, In: Kalkofen W. (Ed.), Numerical Radiative Transfer,
Cambridge, p. 101
\item{}Flower D.R., Launay J.M.  1985, MNRAS 214, 271
\item{}Fuller G.A., Myers P.C. 1992, ApJ 384, 523
\item{}Green S. 1975, ApJ 201, 366
\item{}Green S., Chapman S. 1978, ApJSS 37, 169
\item{}Heaton B.D., Little L.T., Yamashita T., Davies S.R., Cunningham C.T.,
Monteiro T.S. 1993, A\&A 278, 238
\item{}Kr\"ugel E., Chini R. 1994, A\&A 287, 947
\item{}Le Bourlot J. 1991, A\&A 242, 235
\item{}Martin H.M., Sanders D.B., Hills R.E. 1984, MNRAS 208, 35
\item{}Miesch M.S., Bally J. 1994, ApJ 429, 645
\item{}Monteiro T.S. 1985, MNRAS 214, 419
\item{}Press W.H., Flannery B.P., Teukolsky S.A., Vetterling W.T. 1986
Numerical Recipies (Fortran Version), Cambridge
\item{}Turner B.E., Chan K., Green S., Lubowich D.A., ApJ 399, 114
\end{description}

\end{document}
